% Stage 8 restored fragment — 2.7.2 Non-radiating current configurations (archive lines 9366-9569)
\subsection{Non-radiating current configurations}
The section introduction displayed balanced loops, counter-directed pairs, and below-cutoff
guides as laboratory signatures of $\ker\mathcal{R}_\omega$. Subsection~\ref{sec:ch272}
gives the definitive constructive account: admissible $\Jimp^{\mathrm{(NR)}}\in\ker\mathcal{R}_\omega$
with $\mathcal{S}_\omega[\Jimp^{\mathrm{(NR)}}]\neq\mathbf{0}$, developing Eq.
Eq.~\eqref{eq:ch2.nr-current-template} \cite{stratton1941}. Radiation integrals, Huygens surface forms
Eqs.~\eqref{eq:ch2.huygens-equivalent-integral}--\eqref{eq:ch2.huygens-tangential-match}, and
equivalent-source ambiguity Eq.~\eqref{eq:ch2.equivalent-source-ambiguity} from
Subsections~\ref{sec:ch251} and~\ref{sec:ch244} are presupposed; dyadic kernel tabulations from Subsection~\ref{sec:ch231} are not repeated \cite{stratton1941}. Evanescent overlap and
measurement nullity belong to Subsection~\ref{sec:ch273}; worked example examples appear in
Subsection~\ref{sec:ch2102} by reference \cite{harrington2001}.
Assumptions. Fix \(\omega>0\), compact support \(\Omega_s\), admissible impressed data
Eq.~\eqref{eq:ch2.source-admissibility}, charge continuity
Eq.~\eqref{eq:ch1.charge-continuity-harmonic}, outgoing radiation condition
Eqs.~\eqref{eq:ch2.outgoing-radiation-refinement}--\eqref{eq:ch2.outgoing-resolvent-unique},
and null space Eq.~\eqref{eq:ch2.radiation-nullspace-def}
\cite{stratton1941}. All configurations below are required to satisfy
\(\mathcal{D}_{J}[\Jimp^{\mathrm{(NR)}}]=1\) and \(\divg\Jimp^{\mathrm{(NR)}}-\jj\omega\rho=0^{\mathrm{(NR)}}\)
with \(\rho^{\mathrm{(NR)}}\) compatible with boundary class
Eq.~\eqref{eq:ch1.outer-boundary-classes} \cite{stratton1941}. regime qualification
Eq.~\eqref{eq:ch2.regime-component-criterion} governs whether near-zone reactive content may be
reported as transport data \cite{harrington2001}. Non-radiating current (locked definition). A \emph{non-radiating current configuration} is an admissible impressed current
\begin{equation}
\label{eq:ch2.nr-current-def}
\Jimp^{\mathrm{(NR)}}\in\ker\mathcal{R}_\omega,
\qquad
\mathcal{S}_\omega[\Jimp^{\mathrm{(NR)}}]\neq\mathbf{0}\ \text{in general},
\end{equation}
with radiative component \(\mathbf{u}_{\mathrm{rad}}=\mathbf{0}\) and non-radiative component
\(\mathbf{u}_{\mathrm{nr}}=\mathcal{S}_\omega[\Jimp^{\mathrm{(NR)}}]\) from
Eqs.~\eqref{eq:ch2.radiative-component-def}--\eqref{eq:ch2.nonradiative-component-def}
\cite{stratton1941}. Equation~\eqref{eq:ch2.nr-current-def} is the Subsection
~\ref{sec:ch272} development of Eq.~\eqref{eq:ch2.nr-current-template} and
expresses Eq.~\eqref{eq:ch2.null-radiation} with explicit $\mathcal{D}_J=1$: far-zone export
vanishes through Eq.~\eqref{eq:ch2.radiation-integral-outgoing-upgrade}, while near-zone fields
may remain substantial kernel image on field space (Subsection~\ref{sec:ch271})
\cite{stratton1941}. Linearity Eq.~\eqref{eq:ch2.radiation-linearity} implies
admissible superpositions of non-radiating configurations remain non-radiating when charge
continuity is preserved \cite{stratton1941}.
Maxwell-closed constraints. Non-radiating configurations are not arbitrary curl fields. Impressed data must satisfy
\begin{equation}
\label{eq:ch2.nr-current-admissibility-law}
\Jimp^{\mathrm{(NR)}}\in L^{2}(\Omega_s)^{3},
\quad
\divg\Jimp^{\mathrm{(NR)}}-\jj\omega\rho=0^{\mathrm{(NR)}},
\quad
\Jimp^{\mathrm{(NR)}}\ \text{BC-compatible on}\ \Gamma,
\end{equation}
as in Eq.~\eqref{eq:ch2.source-admissibility} \cite{stratton1941}. Equation
\eqref{eq:ch2.nr-current-admissibility-law} excludes discontinuous \(\Jimp^{\mathrm{(NR)}}\) that
violate Eq.~\eqref{eq:ch1.charge-continuity-harmonic} before Maxwell field specification is attempted
\cite{stratton1941}. Surface-equivalent realizations on \(\Gamma_{\mathrm{eq}}\) inherit trace
class constraints from Subsection~\ref{sec:ch211} \cite{harrington2001}. Membership in
\(\ker\mathcal{R}_\omega\) is an \emph{export} statement; Eq.~\eqref{eq:ch2.nr-current-admissibility-law}
is a \emph{source} statement---both are required \cite{stratton1941}.
Let \(\mathcal{C}\subset\Omega_s\) be a closed admissible circulation path enclosing negligible
magnetic flux linkage at \(\omega\), with oppositely directed arms \(\mathcal{C}_+\) and
\(\mathcal{C}_-\) carrying equal magnitude currents. Physically, picture a wire loop in the
\(xy\)-plane with counter-directed segments along \(\pm\hat{\mathbf{x}}\): each arm radiates a
leading \(1/r\) transverse field, but the net electric dipole moment vanishes when the
circulation closes. When far-zone transverse contributions from the two arms cancel under
outgoing propagation,
\begin{equation}
\label{eq:ch2.nr-loop-cancellation}
\int_{\mathcal{C}_+}\!\mathbf{J}_{\mathrm{loop}}\cdot d\boldsymbol{\ell}
\int_{\mathcal{C}_-}\!\mathbf{J}_{\mathrm{loop}}\cdot d\boldsymbol{\ell}
=0
\ \Longrightarrow\
\mathcal{R}_\omega[\Jimp^{(\mathrm{loop})}]=\mathbf{0},
\end{equation}
for the associated volume current \(\Jimp^{(\mathrm{loop})}\) obtained by compact support
regularization on \(\mathcal{C}\) \cite{stratton1941}. Equation
\eqref{eq:ch2.nr-loop-cancellation} is the dipole-moment cancellation law at the source level:
vanishing leading \(1/r\) transverse term in
Eq.~\eqref{eq:ch2.far-field-amplitude-integral} when the net moment vanishes
\cite{jackson1999}. Near-zone reactive fields from \(\mathcal{S}_\omega[\Jimp^{(\mathrm{loop})}]\)
need not vanish; Eq.~\eqref{eq:ch2.regime-component-criterion} still applies \cite{stratton1941}.
Two compact sources \(\Jimp^{(+)}\), \(\Jimp^{(-)}\) with supports \(\Omega_\pm\) separated by
distance \(d\ll\lambda_0\) can form a non-radiating pair when their outgoing far patterns
interfere destructively. A concrete arrangement is a pair of equal-strength electric dipoles
along \(\hat{\mathbf{z}}\), separated by \(d=\lambda_0/20\), driven in opposite phase so the
leading \(1/r\) terms cancel on \(S^{2}\) while reactive overlap remains in the gap:
\begin{equation}
\label{eq:ch2.nr-counterprop-pair}
\Jimp^{\mathrm{(NR,cp)}}=\Jimp^{(+)}+\Jimp^{(-)},
\qquad
\E_{0,+}(\theta,\phi)+\E_{0,-}(\theta,\phi)=\mathbf{0}
\ \text{on}\ S^{2},
\end{equation}
with \(\E_{0,\pm}\) far amplitudes from Eq.~\eqref{eq:ch2.far-field-amplitude-integral} for
each component \cite{stratton1941}. Equation
\eqref{eq:ch2.nr-counterprop-pair} realizes Eq.~\eqref{eq:ch2.nr-current-def} when phase and
polarization of \(\Jimp^{(\pm)}\) are tuned so \(\Pi_{\mathrm{rad}}\mathcal{S}_\omega[\Jimp^{(\pm)}]\)
cancel on the outgoing branch \cite{stratton1941}. The pair may store substantial reactive
energy in the overlap region through Eq.~\eqref{eq:ch2.component-power-budget} while
\(\mathcal{P}_{\mathrm{rad}}=0\) on qualified shells \cite{harrington2001}.
On a closed equivalent surface \(\Gamma_{\mathrm{eq}}\) enclosing radiating sources, tangential
non-radiating distributions satisfy Eq.~\eqref{eq:ch2.equivalent-source-ambiguity}:
\begin{equation}
\label{eq:ch2.nr-surface-tangential-null}
\mathbf{J}_{s,\mathrm{eq}}^{\mathrm{(NR)}}
\in
\ker\mathcal{R}_\omega\big|_{\Gamma_{\mathrm{eq}}},
\qquad
\mathbf{J}_{s,\mathrm{eq}}^{\mathrm{(NR)}}\cdot\hat{\mathbf{n}}=0
\ \text{on}\ \Gamma_{\mathrm{eq}},
\end{equation}
with \(\hat{\mathbf{n}}\) the outward normal \cite{stratton1941}. Equation
\eqref{eq:ch2.nr-surface-tangential-null} is the Subsection~\ref{sec:ch272} constructive realization of
Huygens equivalence ambiguity: exterior fields on \(\mathcal{F}_{\mathrm{SF}}\) are fixed by
exterior radiative uniqueness (Subsection~\ref{sec:ch243}), but equivalent surface currents are unique
only up to additions in Eq.~\eqref{eq:ch2.nr-surface-tangential-null}
\cite{stratton1941}. Integrals Eqs.~\eqref{eq:ch2.huygens-equivalent-integral}--\eqref{eq:ch2.huygens-tangential-match}
realize the same \(\mathcal{S}_\omega\) image for equivalent and non-radiating tangential
perturbations \cite{harrington2001}.
When radiation eigenpairs with \(\lambda_n=0\) are declared in
Eq.~\eqref{eq:ch2.radiation-eigenpair-def},
\begin{equation}
\label{eq:ch2.nr-modal-eigenrealization}
\Jimp^{\mathrm{(NR,mod)}}
=
\sum_{n\in\mathcal{N}_{\mathrm{null}}}
\alpha_n\,\Jimp^{(n)},
\qquad
\mathcal{R}_\omega[\Jimp^{\mathrm{(NR,mod)}}]=\mathbf{0},
\end{equation}
for coefficients \(\alpha_n\in\mathbb{C}\) such that the superposition satisfies
Eq.~\eqref{eq:ch2.nr-current-admissibility-law} \cite{stratton1941}. Equation
\eqref{eq:ch2.nr-modal-eigenrealization} specializes Eq.~\eqref{eq:ch2.modal-null-sector-reduction}
from Subsection~\ref{sec:ch271} to constructive sources: each \(\Jimp^{(n)}\) with
\(\lambda_n=0\) is a modal non-radiating generator when admissible
\cite{stratton1941}. On exterior domains \(\mathcal{N}_{\mathrm{null}}\) need not span all of
\(\ker\mathcal{R}_\omega\) Eq.~\eqref{eq:ch2.nullspace-dimensional-classification}
\cite{harrington2001}.
\begin{figure}[t]
\centering
\ChOneFigBlock{%
\begin{tikzpicture}[ch1fig, node distance=5mm and 7mm]
 \node[ch1box] (loop) {Loop\\cancellation};
 \node[ch1box, right=7mm of loop] (cp) {Counter-\\prop pair};
 \node[ch1box, right=7mm of cp] (surf) {Surface\\tangential};
 \node[ch1box, right=7mm of surf] (mod) {Modal\\$\lambda_n=0$};
 \node[ch1box, below=5mm of cp] (ker) {$\ker\mathcal{R}_\omega$};
 \node[ch1box, below=5mm of loop] (j) {$\Jimp^{\mathrm{(NR)}}$};
 \draw[ch1flow] (loop.east) -- (loop.east |- ker.west);
 \draw[ch1flow] (loop.east |- ker.west) -- (ker.west);
 \draw[ch1flow] (cp.south) -- (cp.south |- ker.north);
 \draw[ch1flow] (cp.south |- ker.north) -- (ker.north);
 \draw[ch1flow] (surf.east) -- (surf.east |- ker.west);
 \draw[ch1flow] (surf.east |- ker.west) -- (ker.west);
 \draw[ch1flow] (mod.east) -- (mod.east |- ker.west);
 \draw[ch1flow] (mod.east |- ker.west) -- (ker.west);
 \draw[ch1flow] (j.south) -- (j.south |- loop.north);
 \draw[ch1flow] (j.south |- loop.north) -- (loop.north);
 \draw[ch1flow] (j.east) -- (j.east |- cp.west);
 \draw[ch1flow] (j.east |- cp.west) -- (cp.west);
 \draw[ch1flow] (j.east) -- (j.east |- surf.west);
 \draw[ch1flow] (j.east |- surf.west) -- (surf.west);
 \draw[ch1flow] (j.east) -- (j.east |- mod.west);
 \draw[ch1flow] (j.east |- mod.west) -- (mod.west);
\end{tikzpicture}%
}
\caption{Subsection~\ref{sec:ch272} canonical non-radiating families: loop cancellation
Eq.~\eqref{eq:ch2.nr-loop-cancellation}, counter-propagating pair
Eq.~\eqref{eq:ch2.nr-counterprop-pair}, surface tangential null
Eq.~\eqref{eq:ch2.nr-surface-tangential-null}, and modal eigenrealization
Eq.~\eqref{eq:ch2.nr-modal-eigenrealization}, all mapping to
\(\ker\mathcal{R}_\omega\).}
\label{fig:ch2.nr-current-configurations}
\end{figure}
The four upper nodes are constructive \emph{families}, not interchangeable templates. Loop cancellation annihilates the leading $1/r$ transverse moment; counter-propagating pairs cancel on $S^{2}$ through phase interference; surface tangential nulls exploit Huygens equivalence ambiguity; modal $\lambda_n=0$ generators realize the spectral slice Eq.~\eqref{eq:ch2.modal-null-sector-reduction}. All four converge on $\ker\mathcal{R}_\omega$, yet each leaves distinct reactive near fields $\mathcal{S}_\omega[\Jimp^{\mathrm{(NR)}}]\neq\mathbf{0}$---the schematic warns that near-zone scans cannot distinguish families once far patterns vanish \cite{stratton1941}. Figure~\ref{fig:ch2.nr-current-configurations} orders constructive families before
evanescent and measurement extensions in Subsection~\ref{sec:ch273}. Equivalent-current solvers on \(\Gamma_{\mathrm{eq}}\) seek \(\mathbf{J}_{\mathrm{eq}}\) matching
near or far data \cite{stratton1941}. Canonical non-radiating configurations (Subsection~\ref{sec:ch272}) and Eq.~\eqref{eq:ch2.nr-surface-tangential-null}
show why such inversions are ill-posed even when forward fields are unique: any solution may
be shifted by tangential \(\mathbf{J}_{s,\mathrm{eq}}^{\mathrm{(NR)}}\) without changing
\(\mathcal{R}_\omega[\Jimp]\) \cite{stratton1941}. Tikhonov regularization
Eq.~\eqref{eq:ch2.tikhonov-source-reconstruction} selects a particular representative but does
not eliminate the equivalence class \cite{harrington2001}. Subsection~\ref{sec:ch273} separates
this source nullity from finite-rank observation nullity
Eq.~\eqref{eq:ch2.measurement-nullspace-composition} \cite{stratton1941}. Subsection~\ref{sec:ch272} assumes linear Maxwell theory at fixed
\(\omega\), outgoing radiation condition, and admissible sources Eq.~\eqref{eq:ch2.nr-current-admissibility-law}
\cite{stratton1941}. Lossy media modify loop and pair conditions through complex
\(k(\mathbf{r})\) without promoting evanescent storage to export
Eq.~\eqref{eq:ch2.rad-null-on-evanescent} \cite{stratton1941}. Canonical families generate
\emph{explicit} members of \(\ker\mathcal{R}_\omega\); they do not exhaust
\(\dim\ker\mathcal{R}_\omega=\infty\) on exterior domains
Eq.~\eqref{eq:ch2.nullspace-dimensional-classification} \cite{harrington2001}.
regime qualification Eq.~\eqref{eq:ch2.regime-operator-screen} remains mandatory: large near-zone
\(|\E|\) from \(\mathcal{S}_\omega[\Jimp^{\mathrm{(NR)}}]\) is not export without
\(\Pi_{\mathrm{rad}}\) qualification \cite{stratton1941}. \phantomsection\label{sec:ch273}
